\section{Methodology}

We design a controlled experiment to compare token entropy and semantic consistency as hallucination predictors, then test whether their combination improves detection.

\subsection{Dataset}

We use \textbf{TriviaQA} \citep{joshi2017triviaqa}, specifically the \texttt{rc.nocontext} split from HuggingFace. This configuration presents questions without accompanying documents, testing the model's parametric knowledge rather than reading comprehension. TriviaQA provides ground-truth answers for each question, enabling automatic evaluation of response correctness.

\textbf{Correctness labeling:} A response is marked correct if it achieves either (1) exact match with any reference answer after normalization, or (2) token-level F1 $>$ 0.5 with the reference.

\subsection{Model}

We use \textbf{Llama-2-7B-chat} \citep{touvron2023llama}, an instruction-tuned variant of the Llama-2 7B base model. This model represents a widely-deployed class of open-source LLMs with full logit access, enabling entropy computation.

\subsection{Response Generation}

For each question, we generate \textbf{10 independent responses} using nucleus sampling with temperature 0.7, top-p 0.95, and max new tokens 128. We use a fixed random seed (42) for reproducibility.

\subsection{Token Entropy Computation}

We compute \textbf{Shannon entropy} from the token-level softmax distributions:

\begin{equation}
H(y_t) = -\sum_{v \in V} p(v | y_{<t}) \log p(v | y_{<t})
\end{equation}

We aggregate by taking the mean across all generated tokens in the response, then average across all 10 responses to obtain a single entropy score per question.

\subsection{Semantic Consistency Computation}

We measure consistency via \textbf{pairwise embedding similarity} across the 10 generated responses:

\begin{enumerate}
    \item Encode each response using SentenceTransformer (\texttt{all-MiniLM-L6-v2})
    \item Compute cosine similarity between all 45 pairs
    \item Average pairwise similarities to obtain the consistency score
\end{enumerate}

Higher consistency indicates stable model behavior; lower consistency indicates hallucination-prone responses.

\subsection{Linear Fusion}

To test combination, we apply \textbf{linear fusion}:

\begin{equation}
\text{Score} = \alpha \cdot \text{Confidence} + \beta \cdot \text{Consistency}
\end{equation}

where Confidence = 1 - normalized\_entropy. Weights $\alpha$, $\beta$ are determined via grid search over [0, 1] with step 0.1 using a held-out validation split.

\subsection{Evaluation Metrics}

\begin{itemize}
    \item \textbf{AUROC:} Area under ROC curve for correctness prediction
    \item \textbf{Cohen's $d$:} Effect size between correct/incorrect distributions
    \item \textbf{Statistical significance:} Two-sample t-test with $\alpha$=0.05
    \item \textbf{Pearson correlation:} Between entropy and consistency scores
\end{itemize}
